\subsection{Twelf Theorems, Proofs, and Wrap-up}
Explaining all of the details of the Twelf code is
beyond the scope of this paper.
So in this section, we conclude the Twelf discussion with a brief
overview of the remaining files of the Twelf $\minilang$ encoding,
how theorems are specified in Twelf, example Twelf proofs,
and how Twelf helps with proving theory about languages.


\subsubsection{Twelf Project Files}
Twelf projects are stored in directories containing a configuration
file typically named \code{sources.cfg}.
The \code{sources.cfg} file tells Twelf the files to read
and the order in which to process them.
The \code{evaluation.elf} file contains the dynamic semantics
of $\minilang$.
File \code{preservation.elf} contains the preservation theorem
and its proof.
File \code{progress.elf} contains the progress theorem
and its proof.
The other files are explained throughout Section~\ref{sec:twelf}.

\subsubsection{Theorems are Function Types. Proofs of Theorems are Function Terms.}
Theorems in Twelf are encoded similarly to
(for all)-(there exists) queries.
For example, below is the encoding of the progress theorem
from lines 11--12 and 153--154 of \code{progress.elf}
(shown without an inline comment).

\begin{verbatim}
progress : of E T -> not_stuck E -> type.
%mode progress +O -NS.
...
%worlds () (progress O NS).
%total O (progress O _).
\end{verbatim}

The first line declares \code{progress} as a type family:
a relation between a derivation of the typing judgment
\code{of E T} and a derivation of the judgment that expression
\code{E} is not stuck (\code{not\_stuck E}).
The \code{\%mode} declaration marks the derivation \code{O}
of \code{of E T} as an input (\code{+}) and the derivation
\code{NS} of \code{not\_stuck E} as an output (\code{-});
it states that for every derivation \code{O}, there exists
a derivation \code{NS}.
The clauses between these declarations, each beginning with
``\code{- :}'', are the cases of the proof.
The \code{\%worlds} declaration states the contexts in which
the theorem holds; the empty parentheses mean that it is only
stated for closed terms, with no hypothetical assumptions.
Finally, \code{\%total} asks Twelf to check that the relation
is \emph{total}:
that the cases cover every possible input derivation
and that the proof terminates, by induction on \code{O}.
If the check succeeds, Twelf has verified the proof.
Twelf also has a more verbose syntax for stating theorems,
the \code{\%theorem} declaration; \code{preservation.elf}
shows the preservation theorem in that syntax in a comment.
For progress, it would be written as follows.

\begin{verbatim}
%theorem
progress :
  forall* {E} {T}
  forall {O:of E T} exists {NS:not_stuck E}
  true.
\end{verbatim}

The \code{forall*} line declares the variables in the theorem.
The next line states that for every derivation of the typing judgment,
\code{of E T}, where the derivation of that judgment is bound to \code{O}, there
exists a derivation of the judgment that expression \code{E} is
not stuck (\code{NS:not\_stuck E}).
The following two rules define the \code{not\_stuck} judgment.\footnote{
In Twelf, function types can be written using arrows
pointing either to the left or to the right.}

\code{not\_stuck/val: not\_stuck E <- value E.} states that if
\code{E} is a value, then \code{E} is not stuck.

\code{not\_stuck/step: not\_stuck E <- step E E'.} states that if
there exists an \code{E'} such that \code{step E E'},
then \code{E} is not stuck as well.

Note that the above progress theorem is nothing more than a function
type.
Given a term/proof of the judgment \code{of E T},
a function of this type should return a term/proof of
the judgment \code{not\_stuck E}.
Hence, a proof of the progress theorem/(function type) is a (total)
function or term of the function type representing the theorem.

\subsubsection{Checking and Deriving Proofs}
The benefit of using Twelf is that it can check your proofs
and occasionally derive proofs.
File \code{test\_typing.elf}, for example, provides two
example judgments that can be derived automatically by Twelf.

\begin{verbatim}
%solve D1 : of (estr (a , b , c , a , eps)) string.
%solve D2 : of (add (enat (s z)) (enat z)) num.
\end{verbatim}

A \code{\%solve} declaration asks Twelf to search for a term
of the given type and, if it finds one, to define a new
constant with that term as its value.
The first declaration asks Twelf to derive that the type of
\code{str[abca]} is \code{str},
and save that derivation in \code{D1}.
The second declaration asks Twelf to derive that the
type of \code{+(num[1]; num[0])} is \code{num},
and save that derivation in \code{D2}.

Below is the Twelf output from those declarations.

\begin{verbatim}
[Opening file test_typing.elf]
%solve
of (estr (a , b , c , a , eps)) string.
 OK
D1 : of (estr (a , b , c , a , eps)) string = of/str.
%solve
of (add (enat (s z)) (enat z)) num.
 OK
D2 : of (add (enat (s z)) (enat z)) num = of/add of/nat of/nat.
[Closing file test_typing.elf]
\end{verbatim}

Twelf finds derivations of both judgments.
It prints them with implicit arguments omitted
(see the footnote in Section~\ref{sec:twelf-static}).
Written with all of their arguments, the derivations are
\begin{verbatim}
D1 = of/str (a , b , c , a , eps).
D2 = of/add (enat z) (enat (s z)) (of/nat z) (of/nat (s z)).
\end{verbatim}
Derivation \code{D1} says that applying function \code{of/str}
to the string term
\code{(a , b , c , a , eps)} returns a
term/derivation of type/judgment
``\code{of (estr (a , b , c , a , eps)) string}''.

Derivation \code{D2} is described as follows:
\begin{enumerate}
\item Apply function \code{of/nat} to term \code{z}
to return a term/derivation of ``\code{of (enat z) num}''.
\item Apply function \code{of/nat} to term ``\code{s z}''
to return a term/derivation of ``\code{of (enat (s z)) num}''.
\item Apply \code{of/add} to the derivations
of ``\code{of (enat z) num}'' and ``\code{of (enat (s z)) num}''
to return a derivation of
``\code{of (add (enat (s z)) (enat z)) num}''.
The first argument \code{(enat z)} establishes
that the third argument should be a term of type
``\code{of (enat z) num}'';
hence, the type of the third argument of \code{of/add}
depends on the first argument passed to \code{of/add}.
The situation is analogous for the second and fourth
arguments of function \code{of/add}.
\end{enumerate}

Twelf could not derive the proof of the preservation theorem,
but it could let you know if your proof was incorrect.
For example, lines 23--27 in \code{preservation.elf}
prove preservation for the typing-evaluation rule combination
(T.4, D.1).
If those lines were removed, Twelf would return the following
error output:

\begin{verbatim}
preservation.elf:98.8-98.11 Error: 
Coverage error --- missing cases:
{E1:exp} {E2:exp} {E3:exp} {O1:of (add E1 E2) num} {S1:step E1 E3}
   {O2:of (add E3 E2) num}
   |- preservation O1 (step/add1 S1) O2.
\end{verbatim}

The error message lets us know that we forgot to prove preservation
for the case when we could derive ``\code{of (add E1 E2) num}''
and ``\code{step (add E1 E2) (add E3 E2)}''
(\code{+($e_1$; $e_2$)} $\stepto$ \code{+($e_3$; $e_2$)}).
In this case, judgment ``\code{step (add E1 E2) (add E3 E2)}''
was derived by applying function (inference rule) \code{step/add1}
to a derivation (\code{S1}) of type ``\code{step E1 E3}''.
To prove preservation for this case,
we would need to produce a derivation of type ``\code{of (add E3 E2) num}''.
This would show that performing an evaluation step on the left
summand of a well-typed \code{add} expression would not change
the type of the resulting \code{add} expression.

In summary, Twelf can aid in deriving simple judgments and checking
that proofs of language properties are correct.
If a proof is not correct, Twelf will let us know that the proof
contains a mistake and provide an error message to help ``debug'' the proof.
